\begin{table}[t]
\tbl{\label{tab:nt}}
{\begin{tcolorbox}[tabularx={p{3cm}|p{2cm}|p{6.3cm}|},enhanced,fonttitle=\bfseries\large,fontupper=\normalsize\sffamily,
colback=yellow!10!white,colframe=red!50!black,colbacktitle=blue!30!white,
coltitle=black,title=Number Theoretic Commands]
\hline
Command & Arguments & Purpose \\
\hline
\hline
\verb'-sift_smooth'  &  $a$ $n$ primes   &  Computes all smooth numbers in the interval $[a,a+n-1]$. Smooth means that they factor completely over the list of primes given. \\
\hline
\verb'-random'  &  $n$ fname  &  Creates $n$ random numbers and writes them to the csv file \verb'fname' \\
\hline
\verb'-random_last'  &  $n$   &  Creates $n$ random numbers prints the last one \\
\hline
\verb'-affine_sequence'  &  $a$ $b$ $p$   &  Splits the interval $[0,p-1]$ into affine sequences of the form $x_{n+1} = a x_n + b \mod p$ \\
\hline
\end{tcolorbox}}
\end{table}
